\documentclass[10pt,a4paper]{amsart}
\usepackage[margin=19mm]{geometry}
\usepackage{amsmath,amssymb,mathtools}
\usepackage[scr=rsfs,cal=euler]{mathalfa}
\usepackage{xfrac}
\usepackage[unicode,hidelinks,bookmarks=false]{hyperref}
\newcommand{\ie}{i.e.\ }
\newcommand{\cf}{cf.\ }
\newcommand{\resp}{respectively\ }
\newcommand{\Subsection}{\subsection}
\providecommand{\nobreakdash}{\nobreak\hbox{-}}
\setlength{\emergencystretch}{2em}
\multlinegap0pt
\usepackage{amssymb}
\usepackage{algorithm}\usepackage{algpseudocode}
\usepackage{xcolor}
\newtheorem{lem}{Lemma}
\newtheorem{thm}{Theorem}
\newcommand{\A}{\mathcal{A}}
\newcommand{\C}{\mathcal{C}}
\newcommand{\La}{\Lambda}
\newcommand{\RR}{\mathbb{R}}
\newcommand{\TT}{\mathbb{T}}
\newcommand{\beq}{\begin{equation}}
\newcommand{\beyy}{\begin{eqnarray*}}
\newcommand{\caI}{\mathcal{I}}
\newcommand{\eeq}{\end{equation}}
\newcommand{\eeyy}{\end{eqnarray*}}
\newcommand{\la}{\lambda}
\newcommand{\p}{\prime}
\newcommand{\tr}{\texttt{tr}}
\title{Tensor Derivatives, Unified Tensor-Form Differential Equations, and Model Reduction via Partial Tucker Decomposition}
\author{Yiran Xu, Changqing Xu}
\date{Source: arXiv:2509.08429v3 (CC BY 4.0)}
\begin{document}
\maketitle

\noindent\emph{Source and licence.} Tensor Derivatives, Unified Tensor-Form Differential Equations, and Model Reduction via Partial Tucker Decomposition. Version arXiv:2509.08429v3 (CC BY 4.0), distributed by arXiv under the Creative Commons Attribution 4.0 International licence, http://creativecommons.org/licenses/by/4.0/, as recorded in the arXiv licence metadata for this exact version and shown on the abstract page https://arxiv.org/abs/2509.08429v3. Full licence notice: \texttt{licenses/arxiv-2509.08429-CC-BY-4.0.txt}.\par\smallskip

\noindent\textit{Editorial scope.} Counted unit: the article's thm th3-2 (source0.tex lines 971--979). The article's complete original proof of it is delivered with it (source0.tex lines 981--1021, one \texttt{proof} environment of 41 lines). No mathematics is written, repaired or re-derived: the statement and the proof are the article's own lines, in the article's own notation, and no retained line is substituted or re-flowed.  Retained uncounted as the source-local context the proof needs: lines 621--627; lines 943--954.  Nothing was omitted from any retained span, so the package carries no omission record.  No retained line contains a citation, so the wrapper carries no bibliography and no bibliography entry.  No retained line contains a figure environment, an image-inclusion command or drawing code, so no image is delivered and the package declares no asset dependency.  Editorial formatting only, in addition: an AMS \texttt{amsart} preamble that restates the article's own declarations and macros used by the retained lines, namely the environments \texttt{lem}, \texttt{thm} on separate counters, together with the standard packages those lines need.  The retained environments print in this document as Lemma 1, Theorem 1 in order of appearance, whereas the article prints the same items with the numbering of its own full document.  The complete native source of the article as distributed in the arXiv e-print for this exact version is bundled unchanged as \texttt{sources/184823/source0.tex}.
\begin{lem}\label{le2-2}
\beq\label{eq:le2-2-1} 
\caI\ast \A = \A = \A\ast \caI, \forall  \A\in\TT_{d;n},
\eeq
where $\ast$ in the first equality is the contractive product of $\caI$ with $\A$ along last $d$ modes of $\caI$, and $\ast$ in the second is the 
contractive product of $\A$ with $\caI$ along the first $d$ modes of $\caI$. 
\end{lem}

\begin{lem}\label{le3-1}
Let $A\in\C^{p\times q}$ be a constant matrix, $X\in\RR^{m\times n}$ a variable matrix, $\la=\la(X)\in\C$ be a function of $X$, and 
$\La=(\la_{ij})\in\C^{m\times n}$ be the matrix with entry $\la_{ij}$ be the derivative of $\la$ w.r.t. $x_{ij}$.  Then 
\begin{description}
\item[(1)]  $\frac{d A}{d X}=0$, i.e., a tensor of size $m\times n\times p\times q$ with all entries being zero.   
\item[(2)]  $\frac{d \la A}{d X}=\La\times A$. 
\item[(3)] $\frac{d(\tr X)}{dX} = I_{n}$ if $m=n$. 
\item[(4)] $\frac{d(\det X)}{dX}= (X^{\ast})^{\top}$ if $m=n$.
\item[(5)] $\frac{d X}{d X}= I_{m}\times_{c} I_{n}$.
\item[(6)] $\frac{d(X^{\top})}{dX}= I_{m}\times_{ac} I_{n}$.
\end{description}
\end{lem}

\begin{thm}\label{th3-2}
Let $A,B$ be constant matrices of appropriate sizes and $Y,Z$ be variable matrices depending on variable matrix $X$. Then we have 
\begin{description}
\item[(a)] $\frac{d (AXB)}{d X}=A^{\top}\times_{c} B$. 
\item[(b)] $\frac{d (YZ)}{d X}=\frac{dY}{dX}\ast_{4} Z +Y\ast_{3}\frac{dZ}{dX}$.
\item[(c)] $\frac{dX^{2}}{d X}=I_{n}\times_{c} X + X^{\top}\times_{c} I_{n}$ if $X\in\RR^{n\times n}$.
\item[(d)]  $\frac{d X^{-1}}{dX} =-X^{-\top}\times_{c} X^{-1}$ if $X$ is invertible.
\end{description}
\end{thm}

\begin{proof}
To prove (a), we may assume that $X\in\RR^{m\times n}$ and $A\in\RR^{p\times m}, B\in\RR^{n\times q}$.  We note that the left hand side (lhs) of (a) is an 4-order tensor with size $m\times n\times p\times q$, which is consistant with that of the right hand side (rhs) of (a). For any 
$(i,j,k,l)\in [m]\times [n]\times [p]\times [q]$, we have 
\beyy 
\left[\frac{d (AXB)}{d X}\right]_{ijkl} 
&=&\frac{d (AXB)_{kl}}{d X_{ij}}\\
&=&\frac{d (\sum_{k^{\p},l^{\p}}A_{kk^{\p}}B_{l^{\p}l}X_{k^{\p}l^{\p}})}{d X_{ij}}\\    
&=&\sum_{k^{\p},l^{\p}}A_{kk^{\p}}B_{l^{\p}l} \frac{d X_{k^{\p}l^{\p}}}{d X_{ij}}\\ 
&=&\sum_{k^{\p},l^{\p}}A_{kk^{\p}}B_{l^{\p}l} \delta_{i k^{\p}}\delta_{j l^{\p}}\\ 
&=&A_{ki}B_{jl} = \left( A^{\top}\times_{c} B\right)_{ijkl}
\eeyy
Thus we have $\frac{d (AXB)}{d X}=A^{\top}\times_{c} B$. \par 
\indent To prove (b), we let $X,Y,Z$ be matrices of size $m\times n, p\times r$ and $r\times q$.  Then the left hand side of (b) is an 4-order tensor of size 
$m\times n\times p\times q$, and it is easy to check that both items in the right hand side of (b) are also 4-order tensors of the same size. Furthermore, 
we have for any $(i,j,k,l)\in [m]\times [n]\times [p]\times [q]$ that  
\beyy 
\left[\frac{d (YZ)}{d X}\right]_{ijkl}&=&\frac{d (YZ)_{kl}}{d X_{ij}}=\frac{d (\sum_{s=1}^{r}Y_{ks}Z_{sl})}{d X_{ij}}\\    
&=&\sum_{s=1}^{r}\left( \frac{d Y_{ks}}{d X_{ij}} Z_{sl} + Y_{ks} \frac{d Z_{sl}}{d X_{ij}}\right) \\ 
&=&\sum_{s=1}^{r}\left[ \left(\frac{d Y}{d X}\right)_{ijks} Z_{sl} + Y_{ks}\left(\frac{d Z}{d X}\right)_{ijsl}\right] \\ 
&=&\left( \frac{dY}{dX}\ast_{4}Z +Y\ast_{3}\frac{dZ}{dX}\right)_{ijsl}
\eeyy 
Thus (b) holds.  Now (c) can be proved by the combination of (b) and (5) of Lemma \ref{le3-1}. In fact, from (b), we have 
 \beyy 
\frac{d (X^{2})}{d X}
& = &X\ast_{3}\frac{dX}{d X} + \frac{dX}{dX}\ast_{4} X \\
& = &X\ast_{3}(I_{n}\times_{c} I_{n}) + (I_{n}\times_{c} I_{n}) \ast_{4} X \\ 
& = &X^{\top}\times_{3} I_{n} + I_{n}\times_{c} X     
\eeyy
The last equality comes from (3) and (4) of  Lemma \ref{le2-2}. \par
\indent To prove (d), we notice that $XX^{-1} = (\det X) I_{n}$.  By (b) and (5) of Lemma \ref{le3-1}, we have 
\beyy
\frac{d(XX^{-1})}{dX} 
&=& \frac{dX}{dX}\ast_{4} X^{-1} + X\ast_{3}\frac{d X^{-1}}{dX}\\
&=& (I_{n}\times_{c} I_{n}) \ast_{4} X^{-1} + X\ast_{3} \frac{d X^{-1}}{dX}\\
&=& I_{n}\times_{c} X^{-1} + X\ast_{3} \frac{d X^{-1}}{dX}
\eeyy
The last equality is due to (4) of Lemma \ref{le2-2}.  Thus we have 
\beq\label{eq:le3-2-1}
X\ast_{3} \frac{d X^{-1}}{dX} = - I_{n}\times_{c} X^{-1} 
\eeq
\end{proof}

\end{document}
